% PAGE-BUDGET TARGET: 6 pages
% HARD MAX: 7 pages
\section{Main theorem, scope boundaries, and exact endpoint}\label{sec:scope}

This paper proves a vacuum-sector mass-gap theorem for the bounded-region local gauge-invariant sharp-local net generated by flowed curvature composites for every compact connected simple gauge group $G$. The mass gap concerns the vacuum Hamiltonian in the OS/Wightman reconstruction of that local net. Faithful-Wilson universality enters only qualitatively: within the admissible gradient-flow renormalization class and after coherent normalization, different faithful ultraviolet Wilson regularizations yield the same continuum local theory, while the explicit Route~1 quantitative certificate remains $\rho$-indexed. Extended-support line and surface sectors remain outside the theorem-level claim. We state those boundaries first because every later reduction in the paper is organized to preserve them exactly.

\subsection{Main theorem and claim boundaries}

\begin{theorem}[Main theorem]\label{thm:main}
Fix a compact connected simple Lie group $G$, choose a faithful finite-dimensional unitary Wilson regularization $\rho:G\to U(d_\rho)$, and fix a weak-window target coupling $u_{\mathrm{ren}}\in(0,u_*]$ for the tuned dyadic construction. Let $\mathcal A_{\mathrm{loc}}(G)$ denote the resulting bounded-region local gauge-invariant sharp-local net generated by flowed curvature composites, with exact theorem-level endpoint clarified in Section~\ref{sec:endpoint}. Throughout the theorem-level group-only endpoint, faithful-Wilson identification is understood within the admissible gradient-flow renormalization class and after the coherent normalization conditions fixed in the endpoint packet. Then the following hold.
\begin{enumerate}[label=\textup{(\Alph*)}]
    \item \textbf{One-shot entrance.} There exist explicit constants and thresholds depending on $(G,\rho)$, including $q^*(G,\rho)\in(0,1)$, a bounded physical entrance constant $\mathrm{Cent}(G,\rho)<\infty$, and an entry index $N_{\mathrm{ent}}(G,\rho;u_{\mathrm{ren}})$, such that the exact one-shot blocked specification enters the KP/Dobrushin basin at a bounded physical scale and remains there along the tuned dyadic trajectory.
    \item \textbf{Dense-core Euclidean clustering.} On a dense positive-time local probe core the Euclidean two-point functions decay exponentially at rate
    \[
      m_*(G,\rho;u_{\mathrm{ren}})=\frac{m_{\mathrm{blk}}(G,\rho)}{2\,\mathrm{Cent}(G,\rho)\,\ell_0},\qquad m_{\mathrm{blk}}(G,\rho):=-\log q^*(G,\rho)>0.
    \]
    \item \textbf{Continuum local gauge-invariant Yang--Mills theory.} There exists a unique continuum sharp-local state on $\mathcal A_{\mathrm{loc}}(G)$ with Euclidean OS strength, Wightman/Haag--Kastler reconstruction, field correspondence for the renormalized gauge-invariant local differential polynomials produced by Lane~A, and a non-triviality witness.
    \item \textbf{Vacuum mass gap.} If $(\mathcal H_{\mathrm{loc}},\Omega_{\mathrm{loc}},H_{\mathrm{loc}})$ denotes the OS/Wightman reconstruction of that sharp-local state, then
    \[
      \operatorname{Spec}(H_{\mathrm{loc}})\subset \{0\}\cup [m_*(G,\rho;u_{\mathrm{ren}}),\infty),\qquad \ker(H_{\mathrm{loc}})=\mathbb C\Omega_{\mathrm{loc}}.
    \]
\end{enumerate}
The continuum local theory depends only on $G$ at theorem level up to the canonical faithful-Wilson identification proved in Section~\ref{sec:endpoint}, within the admissible gradient-flow renormalization class and after coherent normalization. The present explicit Route~1 quantitative certificate remains indexed by the chosen faithful regularization $\rho$.
\end{theorem}
\proofhome{Core Section~8, with Companions~I--III cited theorem-by-theorem}
\proofsynopsis{Part~(A) is assembled from the ultraviolet gate and one-shot entrance module; part~(B) is the downstream Euclidean clustering packet; part~(C) is the Lane~A plus endpoint construction packet; and part~(D) is the Route~1 lattice-side gap followed by the unique QE3 transport chain. Section~8 will give the theorem-chain assembly only, with all line-by-line analytic proofs carried by the companions.}

The theorem-level object is deliberately local. Faithful-Wilson universality identifies the continuum local theory for faithful ultraviolet regularizations of the same compact simple group, but the present manuscript does not upgrade the explicit quantitative lower-bound ledger to a faithful-$\rho$-uniform statement. The endpoint likewise excludes extended-support line and surface sectors from the theorem claim itself.

\nonclaim{The present manuscript does \emph{not} claim a faithful-$\rho$-uniform explicit quantitative entrance or mass-gap bound. The continuum local theory is faithful-Wilson universal, but the explicit Route~1 quantitative certificate remains indexed by the chosen faithful regularization.}

\nonclaim{The present manuscript does \emph{not} claim a vacuum gap for extended-support or other nonlocal sectors. The gap statement concerns only the vacuum Hamiltonian in the OS/Wightman reconstruction of the bounded-region local gauge-invariant sharp-local net.}

\nonclaim{The theorem-level group-only endpoint is asserted only within the admissible gradient-flow renormalization class and after coherent normalization. The present manuscript does \emph{not} claim a scheme-independent identification outside that class or without those normalization conventions.}

\subsection{Effective gauge group and local-net endpoint}

\begin{definition}[Effective gauge group determined by the ultraviolet action]\label{def:effective-gauge-group}
Let $G$ be a compact Lie group and let $\rho:G\to U(d_\rho)$ be the representation used to define the Wilson plaquette weight. Set
\[
  K_\rho:=\ker(\rho)\subset Z(G),\qquad G_{\mathrm{eff}}:=G/K_\rho.
\]
The ultraviolet action, the Wilson lattice measure, and the strictly local gauge-invariant observables considered in this paper depend on $(G,\rho)$ only through $G_{\mathrm{eff}}$.
\end{definition}

\begin{lemma}[Center-blindness and descent to the quotient]\label{lem:center-blindness}
With the Wilson action defined from $\rho$, the one-plaquette weight, the full Wilson lattice gauge measure, and every gauge-invariant observable generated by flowed local curvature composites depend on $G$ only through the quotient $G_{\mathrm{eff}}$. Equivalently, the resulting Schwinger functions and the reconstructed sharp-local OS/Wightman theory are canonically identified for $(G,\rho)$ and $(G_{\mathrm{eff}},\rho_{\mathrm{desc}})$.
\end{lemma}

\begin{proof}
The plaquette action density has the form
\[
  A_\rho(U)=1-\frac{1}{d_\rho}\operatorname{Re}\operatorname{Tr}_\rho(U).
\]
If $k\in K_\rho$, then $\rho(kU)=\rho(U)$ and therefore $A_\rho(kU)=A_\rho(U)$. The Wilson weight is constant on cosets of $K_\rho$, and the lattice measure is invariant under pointwise multiplication of link variables by elements of $K_\rho$. Any gauge-invariant observable built from $\rho$, including flowed local composites of the curvature, is likewise $K_\rho$-invariant, so its expectation depends only on the induced variables in $G/K_\rho$. Passing to the continuum limit and OS reconstruction preserves these identities.
\end{proof}

\begin{lemma}[Adjoint-local generators and insensitivity to global form]\label{lem:adjoint-local}
The sharp-local algebra $\mathcal A_{\mathrm{loc}}(G)$ is generated by renormalized local composite insertions obtained as small-flow-time limits of flowed curvature composites. Every such generator is adjoint-local: it depends on the gauge field only through the associated adjoint bundle and is therefore invariant under central gauge transformations. Consequently, for any central subgroup $K\leq Z(G)$, the local net generated by $\mathcal A_{\mathrm{loc}}$ is canonically identified for $G$ and $G/K$.
\end{lemma}

\begin{proof}
By construction, the flowed probe algebra is generated by gauge-invariant flowed curvature composites built from $G_{\mu\nu}(t,x)$ and its covariant derivatives. Lane~A identifies each sharp-local insertion as a small-flow-time limit of inverse-flow combinations of these probes. Under a gauge transformation $g(x)\in G$, the flowed curvature transforms by the adjoint action, and the same holds for covariant derivatives. Since the defining polynomials are adjoint-invariant, every generator is unchanged under central gauge transformations. Because the adjoint action is identical for $G$ and $G/K$, the resulting strictly local gauge-invariant net is canonically the same for both global forms.
\end{proof}

\subsection{What the paper does not claim}

\begin{corollary}[Scope of the vacuum-gap statement]\label{cor:scope-gap}
The vacuum spectral gap asserted in \cref{thm:main} is a statement about the OS vacuum representation of the bounded-region local gauge-invariant net generated by $\mathcal A_{\mathrm{loc}}(G)$. In particular, it does not assert a gap in other representations or superselection sectors generated by nonlocal line or surface operators.
\end{corollary}

\begin{proof}
This is an immediate consequence of \cref{lem:center-blindness,lem:adjoint-local}. The theorem-level object is the strictly local gauge-invariant sharp-local net. Nonlocal defects and line/surface sectors are logically downstream extensions of that local theory, not part of the theorem-level endpoint proved here.
\end{proof}

\subsection{The three load-bearing theorem modules}

A first-pass audit does not require a linear read of the full technical corpus. The proof burden concentrates in three theorem modules: the ultraviolet/public-scope gate of Section~\ref{sec:uv}; the Route~1 lattice-side gap and weak-window closure of Section~\ref{sec:route1}; and the Lane~A plus QE3 transport chain that closes the continuum sharp-local vacuum gap. The purpose of the next section is to make that reduction explicit at theorem level before any detailed proof reading begins.

With the claim and its boundaries fixed, the remaining task is to expose the proof as a finite theorem object. The next section therefore replaces the impression of a long monograph by an explicit public packet reduction, names the only live Route~1 and Lane~A chains, and isolates the closure nodes that a first-pass referee must keep in view throughout the paper.
